\begin{figure}[H]
\label{schema:offer}
Jeżeli typ pola nie jest podany to jest to string.
\begin{verbatim}
{
  "id": "UUID oferty",
  "name": "Nazwa oferty",
  "vendor": {
    "name": "Nazwa sprzedawcy"
  },
  "price": {
    "amount": "Cena oferty w formacie tekstowym",
    "currency": "Kod waluty"
  },
  "description": "Opis oferty",
  "quantity": "Liczba możliwych sprzedaży oferty -- liczba całkowita",
  "version": "Wersja oferty -- liczba całkowita",
  "participantMode": "NONE", "NAMED" lub "GROUP" -- domyślnie "NONE"
}
\end{verbatim}
\caption{Obiekt Offer}
\end{figure}

Pole \texttt{participantMode} określa, jakich danych uczestników wymaga rezerwacja oferty:
\begin{itemize}
  \item \texttt{NONE} -- bilet nieimienny, dane uczestników nie są przyjmowane,
  \item \texttt{NAMED} -- bilet imienny, każda zarezerwowana sztuka wymaga uczestnika
    z polem \texttt{name},
  \item \texttt{GROUP} -- oferta grupowa, wystarczą dane klienta, dane uczestników nie są przyjmowane.
\end{itemize}
Brak pola oznacza \texttt{NONE}.
Tryb ustawia się przy tworzeniu oferty i nie można go później zmienić.